\myparagraph{\bf Generalize stage}
For each refined query data $(\pc,\assignments)$ produced by the $\sample$ and $\extract$ stages, the $\generalize$ stage transforms it into a generalized representation:
\[
(\pc,\assignments)
\rightarrow
(\pc',\assignments').
\]
where $\pc'$ and $\assignments'$ denote the generalized path condition and assignment, respectively. By abstracting symbolic variable names and concrete values, $\generalize$ allows multiple semantically equivalent query data entries to be represented by a single generalized query data entry. Consequently, a single stored query data can be reused for multiple path conditions across software releases.

Algorithm~\ref{alg:ourtool} invokes $\generalize$ after completing the $\sample$ and $\extract$ stages (line~14). It first decomposes each query data into individual constraints and assignments. For example, given the query data $\adata$:
\[
\adata=(\myset{(x>3)\land(x<10)\land(y>2)},
\myset{x\mapsto4,\;y\mapsto3}),
\]
$\generalize$ first converts it into
\[
\myset{
(x>3),
(x<10),
(y>2),
(x\mapsto4),
(y\mapsto3)
}.
\]
The symbolic variables are then replaced by indexed symbolic variables according to their order of appearance. In this example, $x$ and $y$ are replaced with $\avariable_1$ and $\avariable_2$, respectively. Next, the concrete values are transformed into relative offsets from the minimum value of each symbolic variable. Since the minimum values are $3$ for $\avariable_1$ and $2$ for $\avariable_2$, the resulting generalized representation is:
\[
\scalebox{0.88}{$
\myset{
(\avariable_1>\aconstvalue_1),
(\avariable_1<\aconstvalue_1+7),
(\avariable_2>\aconstvalue_2),
(\avariable_1\mapsto\aconstvalue_1+1),
(\avariable_2\mapsto\aconstvalue_2+1)
}.
$}
\]
Finally, $\generalize$ reconstructs the generalized constraints and assignments into the original query-data format:
\[
\begin{aligned}
\adata' = (&\myset{
(\avariable_1>\aconstvalue_1)
\land
(\avariable_1<\aconstvalue_1+7)
\land
(\avariable_2>\aconstvalue_2)
},\\
&\myset{
\avariable_1\mapsto\aconstvalue_1+1,\;
\avariable_2\mapsto\aconstvalue_2+1
}).
\end{aligned}
\]

The resulting generalized query data is stored in $\prevdata$ and reused during symbolic execution of subsequent software releases (line~15). When the generalized representation of a newly generated query data matches a stored generalized query data, $\ourtool$ reconstructs the corresponding concrete assignment by replacing the indexed variables and relative values with the corresponding symbolic variables and concrete values from the new path condition. Consequently, $\ourtool$ determines feasibility and generates a satisfying assignment without invoking the SMT solver.

For example, suppose the newly generated path condition is
$\myset{(z>12)\land(z<19)\land(w>50)}$.
After applying the same abstraction, it matches the generalized query data $\adata'$. Since the minimum values of $z$ and $w$ are $12$ and $50$, respectively, $\ourtool$ reconstructs the concrete assignment by instantiating
$\myset{\avariable_1\mapsto\aconstvalue_1+1,\;\avariable_2\mapsto\aconstvalue_2+1}$
as
$\myset{
z\mapsto13,\;
w\mapsto51
}$,
thereby determining feasibility without invoking the SMT solver.

\iffalse
\subsection{Generalize Stage} \label{sec:generalize}

With the refined dataset, this stage aims to generalize path conditions and assignments in the query data, so that multiple results are stored as a single query data and solver calls for many path conditions are skipped using a query data. At line~7, $\Generalize$ function first decomposes each path condition into individual constraints, and then unions the path condition and assignment sets of each query data. For example, if a query data $\adata$ is given as follows:
{
\[
\scalebox{0.9}{$
\adata = (\myset{(x > 3) \land (x < 10) \land (y > 2)}, \myset{(x \mapsto 4), (y \mapsto 3)}),
$}
\]
}the $\Generalize$ function transforms $\adata$ as follows:
{
\[
\scalebox{0.9}{$
\myset{(x > 3), (x < 10), (y > 2), (x \mapsto 4), (y \mapsto 3)}.
$}
\]
}For the transformed set, it indexes the variables in their order of appearance. In the example, the variable $x$, which appears first in the constraints, is replaced with $\avariable_{1}$, and $y$ is replaced with $\avariable_{2}$. For the indexed variable $\avariable_{i}$, $\ourtool$ relativizes the values based on the minimum value ($\aminvalue$) among all values for the variable $r_{i}$. Applying this process to the example, with $\aminvalue=3$ for $\avariable_{1}$ and $\aminvalue=2$ for $\avariable_{2}$, returns the set as:
{
\[
\scalebox{0.9}{$
\myset{(\avariable_{1} > \aconstvalue_{1}), (\avariable_{1} < \aconstvalue_{1}+7), (\avariable_{2} > \aconstvalue_{2}), (\avariable_{1} \mapsto \aconstvalue_{1}+1), (\avariable_{2} \mapsto \aconstvalue_{2}+1)}.
$}
\]
}Lastly, $\ourtool$ re-designs the result into the same format as the original query data as:
{
\[
\scalebox{0.85}{$
(\myset{(\avariable_{1} > \aconstvalue_{1}) \land (\avariable_{1} < \aconstvalue_{1}+7) \land (\avariable_{2} > \aconstvalue_{2})}, \myset{\avariable_{1} \mapsto \aconstvalue_{1}+1, \avariable_{2} \mapsto \aconstvalue_{2}+1}).
$}
\]
}

Now, while testing the updated release, if the path condition of a newly explored state matches the loaded data, $\ourtool$ returns feasibility and defines an assignment without invoking the solver. For instance, when the path condition $\myset{(z > 12), (z < 19), (w > 50)}$ is explored, $\ourtool$ matches it against the loaded data and defines $\myset{z \mapsto 13, w \mapsto 51}$ without invoking the solver, using 12 and 50 as the minimum values of $z$ and $w$. In our experiments, this stage allows $\ourtool$ to store query results with 15.1\% lower memory cost on average than the naive approach and reduce the number of attempts to match each query explored in the updated release against entries from previous-version testing by 44.6\%, as discussed in Section~\ref{sec:experiments}. 
\fi